“Jeetje, wat ben je groot!”
En de profeet genoot

van de dwepende griet.
Naar zijn bekkengebied

keek hij toen ondeugend.
Er vast op verheugend

dat dan de discipel
een generend schouwspel

daar zou gaan opvoeren.
Doch over haar toeren

reageerde ze niet.
De aanbiddende griet

lachte slechts verlegen.
Ook dat kwam hij tegen.

De profeet keek terug,
prekend op een zandrug,

net boven het beekdal,
met publiek overal,

zittend rond een klein meer.
Zijn woorden daalden neer

op heide en grasland,
waar water elegant

in muisstille lussen
langzaam stroomde tussen

keien en vliegdennen
zijn weg te verkennen.

Hij begon zijn hoorspel.
“Nee, dit is geen voorstel

maar een voorstelling
voor je ontwikkeling.

De innerlijke weg
is de echte uitweg.”

“Men zegt ga in het lijf,
ik zeg ga in je wijf!

Jij wraakt het als geilig.
Ik zeg het is heilig.

Omarm je natte droom.
Lust is de rauwe stroom

van het leven immers.
Word een van de klimmers,

uit kosmisch verlangen
om moksha te vangen,

op de Mount Meru top.
Daar vindt men bovenop

het symbolisch centrum
van het universum,

alwaar hemel en hel,
in het aardse schaakspel

door lust samenkomen.
Waar wordt aangenomen

dat nimfen al tijden
heiligen verleiden

om zo zeker te zijn
of ze wel verlicht zijn.

De deur naar helderheid
en creativiteit

wordt door lust geopend.
Daar vind men doorlopend

het pad van ontwaking.
Kom je in aanraking

met levensenergie!
Dat is jullie missie!

Word je steeds meer van lust
als een brandstof bewust!

Let dan op je adem
als kalmerend centrum.”

Een verhaal uit India
rept over Arjuna,

een prins en een halfgod,
die in het episch plot,

een boogschietwedstrijd wint,
en zo zijn eega vindt;

de prinses Draupadi.
Hij brengt haar naar Kunti,

zijn moeder, die niet kijkt,
wanneer hij met haar prijkt,

pal na hun binnenkomst.
Dat bleek al snel het stomst

wat ze had kunnen doen.
“Kijk eens”, hoorde de oen:

“wat ik heb meegebracht”.
Ze gaf het geen aandacht

en zo zat Kunti mis
toen ze: “wat het ook is

deel het met je broers”, zei.
Arjuna was niet blij

met wat moeder verzon.
Haar eer of dharma kon

niet worden gebroken,
eenmaal uitgesproken,

hoezeer het haar ook speet.
De gehaaide profeet

gebruikte het verhaal
als smoes voor de moraal

om het goddelijke
niet in menselijke

waarheden te gieten.
Zo mocht hij genieten

van zeven apostelen
zonder te worstelen

met morele kritiek.
De profeet gaf repliek

naar de zeven vrouwen,
die hem als getrouwen

blindelings vereerden
en stiekem begeerden.

Geloof door zijn grieten
in gangbare mythen

werd geenszins geoorloofd.
“Wie van jullie gelooft

dat een worm uit vlees groeit
of een vlieg uit stof vloeit?”

Geen van zijn aanhangers
of zelfs zijn kerkgangers

durfde dit gegeven
eerlijk toe te geven.

“Alle wezens ontstaan
niet abrupt of spontaan.

Het is de foute les,
die Aristoteles

en vele geleerden
hardnekkig beweerden.

Elk leven vraagt zijn bron
- zaad, ei in een cocon -

wat leidt tot bevruchting,
groei en ontwikkeling.

Er is geen Genesis,
die onmiddellijk is.

God verstrekt de wetten,
maar doet niet de zetten,

die het leven scheppen.
Laat je niet steeds neppen

door de nepprofeten
die het beter weten!

Wanen dat een brandgans
uit een schelp en bij kans

ineens wordt geboren,
kan ik niet aanhoren.

Men zoekt daarin bewijs
dat God machtig en wijs

de wereld steeds aanstuurt.
Hij een kind had verstuurd

naar een onbevlekte.
Wat een zotte gekte

naar hem die we prijzen.
God moet zich bewijzen?

Maria was een maagd,
daar ze niet werd geplaagd

door aardse drijfveren.
Zij die vals beweren
dat ze niet bevrucht was
zijn als een stuk kompas.

Ze wijzen steevast in
de verkeerde richting

puur uit onwetendheid.
God schiep verscheidenheid

in zijn natuurwetten.
We moeten opletten

dat we een vreemd wezen
niet als het kwaad lezen.

Zijn de wilde mannen
niet Gods wijze plannen

door ons onbegrepen?
Lijkt niet zonder strepen

een zebra op een paard?
Gelijkt niet onbehaard

een aap op ons als mens?
Uitte God niet zijn wens

voor meer verscheidenheid
ofwel diversiteit

in soorten en in ras
met “zag dat het goed was”?”

Een apostel weende
omdat de vrouw meende

dat de profeet verwees
naar haar leven als wees;

bloedvergoten, verhuisd,
verstoten en verguisd,

omdat ze anders was
dan de rest in haar klas.

Hij veegde haar tranen
en verbood de wanen

dat ze heilig waren.
In die prille jaren

van de menselijkheid,
de nog donkere tijd,

dacht men vaak dat “tranen
een weg naar heil banen.”

“Hou op!”, zei de profeet,
toen men die uitspraak deed.

“Een traan is een vloeistof.
Het voert een afvalstof

af die vrijkomt bij stress.”
Hij vervolgde zijn les

met dat God de ogen
tegen het uitdrogen

alzo wil beschermen.
Want bij het kermen,

knijpt men een oogje dicht.
Hij kuste toen het wicht

op haar fronsend voorhoofd.
“Opdat je mij gelooft!”

Wie was de jonge vrouw,
in die treurige rouw?
